\begin{abstract}
Multi-teacher on-policy distillation (MOPD) combines domain specialists into a single student, but how loss normalization, optimizer dynamics, and vocabulary supervision shape its updates and task performance remains unclear. We study these choices with Qwen3-1.7B and four reinforcement-learned teachers. First, token- and response-level averaging produce distinct gradients but similar final scores under top-64 supervision and Adam; the spread is larger under sampled-token supervision with SGD, while sampled-token Adam endpoints are nearly tied. Second, joint training changes about 97\% of FP32 parameters, whereas only 7.3--8.5\% differ after BF16 rounding. The apparent alignment of teacher-specific updates largely reflects Adam's accumulated moments: subtracting the zero-gradient step reduces their mean cosine from above 0.83 to below 0.01 on domain-specific inputs. Third, the top-64 intersection retains over 99.9\% of student probability on average at the evaluated Qwen checkpoints, yet provides no consistent task advantage over sampled-token supervision. In SmolLM3-3B diagnostics, top-64 retains less student probability on some inputs and reduces gradient error relative to top-4 and top-8, but does not reproduce the full-vocabulary gradient. These results distinguish current gradients from optimizer history and stored-weight changes, and show why vocabulary supervision should be assessed through per-task learning curves alongside parameter statistics.
\end{abstract}
